\begin{thebibliography}{99}

\bibitem{aho-hopcroft-ullman}
А.\,Ахо, Дж.\,Ульман, Дж.\,Хопкрофт. \textit{Построение и анализ
вычислительных алгоритмов}, пер. с англ. А.\,О.\,Слисенко
под редакцией Ю.\,В.\,Мати\-ясевича. М.:\,Мир, 1979.

\bibitem{boolos-jeffrey}
        %
Дж.\,Булос, Р.\,Джеффри. \textit{Вычислимость и логика},
пер. с англ. В.\,А.\,Душского и Е.\,Ю.\,Но\-ги\-ной под
редакцией С.\,Н.\,Артёмова.
М.:\,Мир, 1994.\ 396~с.

% COMPUTABILITY AND LOGIC
% Third edition
% George S. Boolos
% Professor of Philosophy
% Massachusetts Institute of Technology
% Richard C. Jeffrey
% Professor of Philosophy
% Princeton University
% Cambridge University Press
% Cambridge / New York / Port Chester / Melbourne / Sydney

\bibitem{bourbaki}

Н.\,Бурбаки.\ \ \textit{Начала математики.\  \ Первая часть.\ \ Основные
структуры анализа. Книга первая. Теория множеств},
пер. с французского Г.\,Н.\,По\-ва\-ро\-ва и
Ю.\,А.\,Шихановича под редакцией В.\,А.\,Ус\-пен\-ско\-го.
М.:\,Мир, 1965.

% \'El\'ements de math\'ematique par N.\,Bourbaki
% XVII, XX, XXII, I
% Premiere partie
% Les structures fondamentales de l'analyse
% livre I
% TH\'EORIE DES ENSEMBLES
% Chapitre I. Description de la math\'ematique formelle
% Chapritre II. Th\'eorie des ensembles. Deuxi\`eme \'edition 1960
%
% Chapitre III. Ensembles ordonn'es -- Cardinaux -- Nombres entiers
% 1956
%
% Chapitre IV. Structures
% 1957
%
% Fascicules de r'esultats, Troisi\`eme \'edition, 1958
%
%  HERMANN

\bibitem{vanderwaerden}

Б.\,Л.\, ван дер Варден. \textit{Алгебра}, перевод с немецкого
А.\,А.\,Бель\-с\-ко\-го. Под редакцией Ю.\,И.\,Мерзлякова. М.:\,Наука,
главная редакция физико\д ма\-те\-ма\-ти\-чес\-кой литературы, 1976.

%  B.L. van der Waerden
%   Algebra I
%  Achte Auflage der Modernen Algebra
%
%  Springer-Verlag
%  Berlin - Heidelberg - New York
%  1971
%
%  Algebra II
%  F\"unfte Auflage.
%
%  Springer-Verlag
%  Berlin - Heidelberg - New York
%  1967

\bibitem{computable-functions}

Н.\,К.\,Верещагин, А.\,Шень. \textit{Лекции по математичес\-кой логике и
теории алгоритмов. Часть 3. Вычислимые фун\-кции}. 3\д е изд.
М.:\,МЦНМО, 2008.
176~с.

\bibitem{naive-set-theory}

Н.\,К.\,Верещагин, А.\,Шень. \textit{Лекции по математичес\-кой логике и
теории алгоритмов. Часть 1. Начала теории множеств}. 3\д изд.
М.:\,МЦНМО, 2008. 128~с.

\bibitem{heyting}

А.\,Гейтинг. \textit{Интуиционизм. Введение}, пер. с
англ. В.\,А.\,Янкова под редакцией и с комментариями
А.\,А.\,Маркова. М.:\,Мир, 1965. 200~с.

% Intuitionism. An introduction.
% A.Heyting
% Professor of Mathematics
% University of Amsterdam
% 1956
% North-Holland Publishing Company
% Amsterdam

\bibitem{hilbert-bernays}

Д.\,Гильберт, П.\,Бернайс. \textit{Основания математики.
Логические исчисления и формализация арифметики}, перевод с
немецкого Н.\,М.\,Нагорного под редакцией С.\,И.\,Адяна.
М.:\,Наука, главная редакция физико\д математической литературы,
1979. 560~с.

% D.\,Hilbert und P.\,Bernays
% Grundlagen der Mathematik. I
% Zweite Auflage. Springer-Verlag
% Berlin -- Heidelberg -- New york
% 1968

\bibitem{gindikin}
        %
С.\,Г.\,Гиндикин. \textit{Алгебра логики в задачах}. М.:\,Наука,
главная редакция физико\д математической литературы, 1972.
288~с.

\bibitem{gladky-textbook}

А.\,В.\,Гладкий. \textit{Математическая логика}. М.:\,Рос\-сий\-с\-кий
государственный гуманитарный университет, 1998. 479~с.

\bibitem{davis-nonstandard}

М.\,Дэвис. \textit{Прикладной нестандартный анализ}, перевод с
англ. С.\,Ф.\,Сопрунова под редакцией и с предисловием
В.\,А.\,Ус\-пен\-с\-ко\-го. М.:\,Мир, 1980. 236~с.

% Applied Nonstandard Analysis
% Martin Davis
% Courant Institute of Mathematical Sciences
% New York University
%
% A Wiley-Interscience Publication
% John Wiley & Sons
% New York - London - Sydney - Toronto
% 1977

\bibitem{ershov-palyutin}

Ю.\,Л.\,Ершов, Е.\,А.\,Палютин. \textit{Математическая логика}. 2-е изд.
М.:\,Нау\-ка, главная редакция физико\д математической литературы,
1987. 336~с. 

\bibitem{keisler-chang}

Г.\,Кейслер, Ч.\,Ч.\,Чэн. \textit{Теория моделей},
перевод с англ. С.\,С.\,Гон\-ча\-ро\-ва, С.\,Д.\,Денисова,
В.\,А.\,Душского и Д.\,И.\,Свириденко. Под редакцией
Ю.\,Л.\,Ершова и А.\,Д.\,Тайманова. М.:\,Мир, 1977. 614~с.

%  C.C.Chang
%       University of California, Los Angeles
%  H.J.Keisler
%       University of Wisconsin, Madison
%  Model  Theory
%
%  (Studies in Logic and the Foundations of Mathematics,
%     volume 73)
%
%  North-Holland Publishing Company, Amsterdam - London
%  American Elsevier Publishing Company, Inc. - New York
%  1973
\bibitem{kurosh}

А.\,Г.\,Курош. \textit{Лекции по общей алгебре}, 2\д е изд.
М.:\,Нау\-ка, главная редакция физико\д математической литературы,
1973. 399~с.

\bibitem{kleene-metamathematics}
С.\,К.\,Клини. \textit{Введение в метаматематику}, перевод с английского
А.\,С.\,Есенина\д Вольпина под редакцией В.\,А.\,Успенского. М.:\,Издательство
иностранной литературы, 1957. 526~с.

%  Introduction to metamathematics
%  Stephen Cole Kleene
%  Profeccor of Mathematics at
%  the University of Wisconsin, Madison
%  D. van Nostrand Company, Inc. New York -- Toronto 1952

\bibitem{kleene-logic}
С.\,К.\,Клини. \textit{Математическая логика}, перевод с английского
Ю.\,А.\,Гастева под редакцией Г.\,Е.\,Минца. М.:\,Мир, 1973. 480~с.

%  Mathematical Logic
%  Stephen Cole Kleene
%  Cyrus C. MacDuffee Profeccor Of Mathematics
%  The University of Wisconsin, Madison
% John Wiley & Sons, Inc. New York -- London -- Sydney, 1967

\bibitem{kleene-vesley}
С.\,Клини, Р.\,Весли. \textit{Основания интуиционистской
математики с точки зрения теории рекурсивных функций}, перевод с
английского Ф.\,А.\,Кабакова и Б.\,А.\,Кушнера.
М.:\,Наука, главная редакция физико\д математической литературы,
1978. 272~с. (Серия: Математическая логика и основания математики.)

%  The Foundations of Intuitionistic Mathematics
%   especially in relation to recursive functions
%  S.C.Kleene  R.E.Vesley
%  North-Holland Publishing Company, Amsterdam, 1965

\bibitem{cormen-leiserson-rivest}
Т.\,Кормен, Ч.\,Лейзерсон, Р.\,Ривест. \textit{Алгоритмы:
построение и анализ}, пер. с англ. К.\,Белова,
Ю.\,Бо\-рав\-лё\-ва, Д.\,Ботина, В.\,Го\-ре\-ли\-ка, Д.\,Дерягина,
Ю.\,Кал\-ниш\-кана, А.\,Ка\-та\-но\-вой, С.\,Львовс\-ко\-го,
А.\,Ро\-ма\-щен\-ко, К.\,Сонина, К.\,Трушкина, М.\,Ушакова,
А.\,Ше\-ня, В.\,Шувалова, М.\,Юдашкина под ред. А.\,Шеня,
В.\,Ященко. М.:\,МЦНМО, 1999. 960~с.

% Thomas H. Cormen, Charles E. Leiserson, Ronald L. Rivest
%  Introduction to Algorithms
%  The MIT Press (Cambridge, Massachusets / London, England)
%  McGraw-Hill Book Company
%    (New York / St.Louis / San Francisco / Montreal / Toronto)
%  1990

\bibitem{lavrov-maksimova}
И.\,А.\,Лавров, Л.\,Л.\,Максимова. \textit{Задачи по теории
множеств, математической логике и теории алгоритмов}, издание
второе. М.:\,Наука, 1984. 224~с.

\bibitem{lyndon}

Р.\,Линдон. \textit{Заметки по логике}, пер. с английского
Ю.\,А.\,Гас\-те\-ва под редакцией И.\,М.\,Яглома. М.:\,Мир, 1968.
128~с.

%  Notes on Logic
%       by
%  Roger C. Lyndon
%  The University of Michigan
%  D. van Nostrand Company, Inc.
%  Princeton, New Jersey
%  Toronto    New York   London
%  1966

\bibitem{manin}
Ю.\,И.\,Манин. \textit{Доказуемое и недоказуемое}.
М.:\,Со\-вет\-с\-кое радио, 1979. 168~с.

\bibitem{robinson}

А.\,Робинсон. \textit{Введение в теорию моделей и метаматематику
алгебры}, пер. с англ. А.\,Б.\,Волынского под редакцией
А.\,Д.\,Тай\-ма\-но\-ва. М.:\,Нау\-ка, главная редакция физико\д
математической литературы, 1967. 376~с. (Серия: Математическая
логика и основания математики.)

%  Abraham Robinson
%   University California, Los Angeles
%  Introduction to Model Tgeory and to the Metamathematics of Algebra
%  North-Holland Publishing Company
%  Amsterdam, 1963

\bibitem{smullyan-informal}
Рэймонд М.\,Смаллиан. \textit{Как же называется эта книга?},
пер. с англ. Ю.\,А.\,Данилова. М.:\,Мир, 1981. 240~с.

% Raymond M. Smullyan
% WHAT IS THE NAME OF THIS BOOK?
% Prentice-Hall, Inc.
% Englewood Cliffs, New Jersey, 1978

\bibitem{smullyan-formal}
Р.\,Смальян.
\textit{Теория формальных систем}, перевод с английского
Н.\,К.\,Ко\-совс\-ко\-го под редакцией Н.\,А.\,Ша\-ни\-на. М.:\,Наука,
главная редакция физико\д математической литературы, 1981.
207~с. (Серия: Математическая логика и основания математики.)

%  THEORY OF FORMAL SYSTEMS,
%  Raymond M. Smullyan
%  Revised Edition
%  Princeton, New Jersey
%  Princeton University Press, 1962

\bibitem{handbook-sets}
\textit{Справочная книга по математической логике в четырёх частях
под ред. Дж.\,Барвайса. Часть~II. Теория множеств},
пер. с англ. В.\,Г.\,Кановея под редакцией
В.\,Н.\,Гришина. М.:\,Наука, 1982. 376~с.

\bibitem{handbook-recursion}
        %
\textit{Справочная книга по математической логике в четырёх частях под
ред. Дж.\,Барвайса. Часть~III. Теория рекурсии},
пер. с английского  С.\,Г.\,Дворникова, И.\,А.\,Лав\-ро\-ва.
Под ред. Ю.\,Л.\,Ер\-шо\-ва.
М.:\,Наука, 1982.\ 360~с.

\bibitem{uspensky-nonstandard}

В.\,А.\,Успенский. \textit{Что такое нестандартный анализ?}
М.:\,Нау\-ка, главн. ред. физико\д математической литературы,
1987. 128~с.

\bibitem{uspensky-nonstandard-small}

В.\,А.\,Успенский.\ \ \textit{Нестандартный, \ или неархимедов,\ \
анализ}. \ М.:\,Знание, 1983. 61~с. (Новое в жизни, науке, технике.
Математика, кибернетика, \номер\,8.)

\bibitem{freudenthal}

Х.\,Фрейденталь. \textit{Язык логики}, перевод с английского Ю.\,А.\,Пет\-ро\-ва
под редакцией Ю.\,А.\,Гастева. М.:\,Наука, главная редакция физико\д
математичес\-кой литературы, 1969. 136~с.

% The Language of Logic
% Hans Freudenthal
% Professor of Pure and Applied Mathematics
% Utrecht University
% Utrecht, The Netherlands
% Elsevier Publishing Company
% Amsterdam -- London -- New York
% 1966

\bibitem{church}

А.\,Чёрч. \textit{Введение в математическую логику. I},
перевод с английского В.\,С.\,Чернявского под ред.
В.\,А.\,Успенского. М.:\,Из\-да\-тель\-с\-тво иностранной литературы,
1960. 484~с.

% Introduction to Mathematical Logic
%    by Alonzo Church
%     volume I
%  Princeton, New Jersey
%  Princeton University Press
%  1956

\bibitem{shoenfield}
Дж.\,Шенфилд. \textit{Математическая логика},
перевод с английского И.\,А.\,Лаврова и И.\,А.\,Мальцева
под редакцией Ю.\,Л.\,Ершова. М.:\,Наука, 1975. 528~с.

\bibitem{engeler}
Э.\,Энгелер. \textit{Метаматематика элементарной математики},
перевод с немецкого Г.\,Е.\,Минца под редакцией А.\,О.\,Слисенко.
М:\,Мир, 1987. 128~с.

%  E.\,Engeler. Metamathematik der Elementarmathematik. Mit 29 Abbildungen
%  (Hochshultext)
%  Springer-Verlag
%  Berlin Heidelberg New York 1983

\bibitem{cyber}
        %
С.\,В.\,Яблонский. \textit{Введение в дискретную матема\-тику},
издание второе. М.:\,Наука, 1986. 384~с.

\bibitem{keisler-nonstandard}

H.\,J.\,Keisler. \textit{Elementary Calculus}.
Weber and Sсhmidt, Prindle, 1976.

\end{thebibliography}
